\section{Implementation: the \texttt{mnlearn} library}
\label{sec:method-implementation}

The system is implemented as a Python package called
\texttt{mnlearn}. The package separates concerns into seven
self-contained subsystems, each exposing a stable public API
through its \texttt{\_\_init\_\_.py}. This section describes
the library layout (Section~\ref{subsec:method-lib-layout}),
the YAML configuration system that drives every experiment
(Section~\ref{subsec:method-yaml}), and the reproducibility
guarantees the runners provide
(Section~\ref{subsec:method-reproducibility}).

\subsection{Library layout}
\label{subsec:method-lib-layout}

The package is organized into seven submodules covering data,
models, inference, learning, baselines, configuration, and
experiment aggregation. A summary of which module owns what is
in \texttt{docs/architecture.md} alongside the source. The
complete source, configs, notebooks, and experiment artifacts
are available at
\url{https://github.com/TomasMajer12/Master-thesis}.

Two top-level entry points run experiments end-to-end. The
structured-prediction path is
\texttt{mnlearn.learning.train(cfg)}, which loads the config,
sets the seed, builds data + model + trainer, fits, evaluates
on the test split, and dumps artifacts. The two-stage baseline
path is \texttt{mnlearn.baselines.train\_baseline(cfg)} and
follows the same API. Both runners share the runtime helpers
in \texttt{mnlearn.learning.runtime}, so the artifact format is
identical across runs.

\subsection{YAML configuration system}
\label{subsec:method-yaml}

Every experiment is driven by a YAML config that maps to a
dataclass schema in \texttt{mnlearn.config}. Three flavors of
YAML exist: an architecture YAML (backbone + graph), an
M\textsuperscript{3}N experiment YAML, and a baseline
experiment YAML; experiment YAMLs reference an architecture
YAML either by relative path or as an inline dict, and both
forms round-trip through the loader. A validator catches
malformed configurations at load time (e.g.\ an OCR baseline
architecture that ends in Softmax, or an M\textsuperscript{3}N
experiment that pairs \texttt{m3n\_hinge} with an LP-relaxation
oracle). A complete reference for every YAML field is
maintained in \texttt{docs/yaml\_reference.md} alongside the
source.

\subsection{Reproducibility}
\label{subsec:method-reproducibility}

Each runner sets every random source the project touches before
constructing the data, model, or trainer. The
\texttt{set\_seed(s)} helper in
\texttt{mnlearn.learning.runtime} seeds Python's
\texttt{random}, NumPy's legacy global RNG (used by the data
builders), PyTorch's CPU RNG, and PyTorch's CUDA RNG on every
device, then sets
\texttt{torch.backends.cudnn.deterministic = True} and
\texttt{cudnn.benchmark = False} so convolutions pick the same
algorithm every run.

Every experiment dumps a fixed artifact set into the resolved
output directory (default \texttt{results/<experiment.name>}):
\begin{itemize}
    \item \texttt{config.yaml} --- the configuration as
    actually executed, with the resolved \texttt{output\_dir}
    baked in. This is a round-trippable input to the loader,
    so a run can be replayed exactly from its own artifact.
    \item \texttt{history.json} --- per-epoch training and
    validation metrics plus run diagnostics ($\phi$ norm for
    LP-M3N, learning rate, etc.).
    \item \texttt{results.json} --- final test-split metrics:
    Hamming and $0/1$ for the structured runner; $0/1$, solver
    feasibility, and OCR clue error for the two-stage baseline
    (Section~\ref{subsec:method-baseline-two-stage}).
    \item \texttt{model.pt} --- the trained
    \texttt{state\_dict()}.
\end{itemize}

Together, the YAML round-trip and the artifact set mean every
result reported in Chapter~\ref{chap:experiments} can be
reproduced from a single \texttt{config.yaml} file plus the
data benchmarks of Chapter~\ref{chap:datasets}.
